\scanpage{131}
De plus, procédant par récurrence sur \(d+d'\), et utilisant une
\uline{factorisation d'Artin} de \(U\)
[ce qui est loisible, car l'on peut prendre \(U\) petit],
on est ramené au cas où \(U\) est de \uline{dim 1}.

Alors la suite exacte pour \(\deleted{\unsure{g}}fg\),
\deleted{\unsure{se}}
compte tenu de
\[
R^0_!f(R^q_!\unsure{g_*|g}(Q_l))=0,
\]
se réduit à \deleted{des} suites exactes courtes
\[
\begin{array}{l}
0\mathbin{\text{---}} R^2_!f_*(R^{n-2}g_*(Q_l))
\mathbin{\text{---}}\unsure{R^n(gf)_!|R^n_!(fg)_*}(Q_l)
\mathbin{\text{---}}R^1_!f_*(R^{n-1}g_*(Q_l))\\
\hfill\mathbin{|}\\[-4pt]
\hfill0
\end{array}
\]
donc par lemme 1 on est ramené à vérifier que
\[
R^n_!\overbrace{\unsure{gf|fg}}^{h}\unsure{_{!*}|_*}(Q_l)
\]
est spécial.
On utilise la résolution des singularités pour \(V\),
on a \(V=\widehat V-Z\), \(\widehat V\) projectif et lisse sur \(k\), d'où
\[
R^{n-1}_!h_Z(Q_l)\longrightarrow
R^n_!\unsure{\widehat h|h}(Q_l\longrightarrow
R^n_!\widehat h(Q_l)
\]
% archive foliation: 130
